\documentclass[11pt]{amsart}

% Encoding and fonts
\usepackage[T1]{fontenc}
\usepackage[utf8]{inputenc}
\usepackage{lmodern}

% Math packages
\usepackage{amsmath}
\usepackage{amssymb}
\usepackage{amsfonts}
\usepackage{amsthm}
\usepackage{mathtools}

% Layout and utilities
\usepackage{geometry}
\usepackage{enumitem}
\usepackage[hidelinks]{hyperref}

\geometry{margin=1in}

\numberwithin{equation}{section}
\setcounter{tocdepth}{3}
\setcounter{secnumdepth}{3}

% ---------------------------------------------------------
% Theorem-like environments
% ---------------------------------------------------------

\theoremstyle{plain}
\newtheorem{theorem}{Theorem}[section]
\newtheorem{proposition}[theorem]{Proposition}
\newtheorem{lemma}[theorem]{Lemma}
\newtheorem{corollary}[theorem]{Corollary}
\newtheorem{conjecture}[theorem]{Conjecture}

\theoremstyle{definition}
\newtheorem{definition}[theorem]{Definition}
\newtheorem{experiment}[theorem]{Experiment}
\newtheorem{heuristic}[theorem]{Heuristic}
\newtheorem{question}[theorem]{Question}

\theoremstyle{remark}
\newtheorem{remark}[theorem]{Remark}
\newtheorem{warning}[theorem]{Warning}

\DeclareMathOperator{\tr}{tr}
\DeclareMathOperator{\rank}{rank}
\DeclareMathOperator{\Var}{Var}
\DeclareMathOperator{\Cov}{Cov}

\title{
Orientation-Bit Geometry of the M\"obius Path:
Reflection Fibres, a Divisor-Defect Gram Form,
and the Recovery of Sine-Kernel Content
from Measured Quantities
}

\author{}
\date{Version 2.0}

\begin{document}

\begin{abstract}
This notebook develops a single geometric object, the Gram form of the
divisor defect on the orientation bits of the M\"obius path, and uses
it to separate what the quadratic profile of that path can and cannot
see.

Three statements are proved here. First, the fibre of the quadratic
profile $n \mapsto M(n)^2$ is exactly an orbit of an elementary abelian
reflection group generated by excursion reflections. Second, the
induced null on excursion signs is exactly isotropic, so that the
expected energy of every eigenspace equals its dimension with no
correction term. Third, and centrally, every cross-excursion
contribution to the quadratic energy vanishes identically,
\begin{equation*}
\sum_{a \in E_j} \sum_{b \in E_k}
\mu(a)\mu(b)\left(N+1-\max(a,b)\right) = 0
\qquad (j \neq k),
\end{equation*}
because the increments of a complete excursion sum to zero. It follows
that no statistic factoring through the quadratic profile can detect
cross-excursion correlation of $\mu$. The Gram form on orientation bits
is introduced precisely because it is reflection-odd and therefore
lives in the complement of that blind spot.

The notebook also isolates a methodological point. A statistic exhibits
sine-kernel behaviour non-tautologically only when its value under the
null is a number different from its value under the sine kernel;
statistics whose null already equals the sine-kernel value, such as the
mean square of $M(x)/\sqrt{x}$ and the energy deficit derived from it,
carry no information about the zeros. Three measured quantities from
separate work do satisfy the criterion, and for each the sine-kernel
coefficient is written in closed form as a function of the measurement:
through the consecutive spacing ratio and the superposition family
$Y_2^{(p)}(x) = p^{-1} Y_2(x/p)$; through the moments
$(n+1)^{-1}(n+2)^{-1}$ of the triangular form-factor deviation, which
determine the sine kernel uniquely as a Hausdorff moment problem; and
through the spacing variance of the lobe sequence.

A list of results established in separate work and used here as inputs
is given in Section~\ref{sec:standing_results}, with the status of each
stated explicitly. Conjectures, heuristics and negative results are
labelled throughout. No numerical computation is performed.
\end{abstract}

\maketitle

\tableofcontents

\bigskip

\begin{remark}
\label{rem:exploratory_nature}

This notebook is exploratory in nature.

Whenever a statement is rigorous it is stated inside a theorem, lemma,
proposition or corollary environment and a proof is supplied. Results
established in separate work are collected in
Section~\ref{sec:standing_results} and are stated without proof, marked
as such. All remaining claims are heuristic or conjectural and are
labelled accordingly.

No numerical computation is performed anywhere in this notebook. The
empirical quantities appearing in
Section~\ref{sec:sine_kernel_from_measurements} are inputs recorded in
separate work, not results obtained here.

\end{remark}

\begin{warning}
\label{warn:bibliographic_metadata}

Bibliographic metadata in Section~\ref{sec:bibliography} has not been
verified against primary sources in this version. DOIs are omitted
wherever they were not confirmed. Entries should be checked before this
notebook is used as a source.

\end{warning}

\section{Introduction}
\label{sec:introduction}

\subsection{The object}
\label{subsec:the_object}

Let $\mu$ denote the M\"obius function and
\begin{equation}
M(x) = \sum_{n \le x} \mu(n)
\label{eq:mertens_definition}
\end{equation}
the Mertens function, regarded as a $\pm 1$ walk supported on the
squarefree integers. We study this walk through three successive
conditionings: on the squarefree mask, on the number of signs of each
type, and on the quadratic profile $n \mapsto M(n)^2$.

The third conditioning is the subject of the notebook. We prove in
Section~\ref{sec:reflection_fibre} that the fibre of the quadratic
profile is exactly an orbit of an elementary abelian reflection group
generated by excursion reflections, and that the uniform measure on
that group is a null model with no free parameter whatsoever. We then
carry a symmetric bilinear form built from the divisor defect down to
the quotient, obtaining a form on the orientation bits of the excursion
decomposition.

\subsection{What is proved here}
\label{subsec:what_is_proved_here}

Three results are proved in this notebook.

\begin{enumerate}[label=(\roman*)]
\item \emph{Reflection fibre}
(Proposition~\ref{prop:reflection_fibre}). The fibre of the quadratic
profile through a sign vector $\varepsilon$ is exactly the orbit
$\mathcal{R}_N \cdot \varepsilon$ of the excursion reflection group
$\mathcal{R}_N \cong (\mathbb{Z}/2)^{Z(N)}$, where $Z(N)$ is the number
of return points of $M$ below $N$. Both inclusions are proved.

\item \emph{Exact isotropy}
(Proposition~\ref{prop:exact_isotropy}). Under the uniform measure
$\nu$ on $\mathcal{R}_N$ the excursion signs satisfy
$\mathbb{E}_\nu[s s^{\mathsf{T}}] = I_Z$ exactly, so that
$\mathbb{E}_\nu\langle s, \widehat{P} s\rangle = \rank \widehat{P}$ with
no correction term, and the variance is given in closed form.

\item \emph{Blindness of the quadratic profile}
(Theorem~\ref{thm:cross_excursion_vanishing} and
Corollary~\ref{cor:profile_blindness}). For $j \neq k$,
\begin{equation}
\sum_{a \in E_j} \sum_{b \in E_k}
\mu(a)\mu(b)\left(N+1-\max(a,b)\right) = 0,
\label{eq:blindness_in_introduction}
\end{equation}
because the sum of the increments over a complete excursion vanishes.
Consequently the quadratic energy depends only on within-excursion
two-point data, and no statistic factoring through the quadratic
profile can detect cross-excursion correlation of $\mu$.
\end{enumerate}

Statement (iii) is the reason the Gram form on orientation bits is
introduced at all: by Proposition~\ref{prop:pushed_form_odd} that form
is reflection-odd, hence lives exactly in the complement of the blind
spot identified by \eqref{eq:blindness_in_introduction}.

\subsection{What is taken as input}
\label{subsec:what_is_taken_as_input}

The notebook draws on a body of results established in separate work.
They are listed in full, with status, in
Section~\ref{sec:standing_results}; two of them are used actively and
the remainder provide context. For orientation, the actively used ones
are:

\begin{itemize}
\item the tensor factorisation
$\widehat{K}(r,q) = \prod_p L_p(v_p(r), v_p(q))$ over all of
$\mathbb{N} \times \mathbb{N}$, with explicit tridiagonal local
factors, which enters as hypothesis (G2) of
Definition~\ref{def:hypotheses_on_G} and is what makes the eigenbasis of
the Gram form a character basis of the multiplicative structure;

\item the exact permutation-null mean
$\mathbb{E}_\pi[E_N]/N = 1/\pi^2 + M(N)^2/(3N) + o(1)$, which is the
second rung of the conditioning hierarchy of
Section~\ref{sec:three_conditionings}.
\end{itemize}

Four further ingredients, referred to as \textbf{(N1)}--\textbf{(N4)},
supply the null model, its error bars, the central construction, and an
interface to Robin-type criteria; they are introduced where they are
used.

\subsection{The methodological point}
\label{subsec:methodological_point}

Section~\ref{sec:sine_kernel_from_measurements} isolates a criterion
that we regard as the main conceptual content of Version 2.0.

The sine kernel is the kernel of a flat band projector
(Lemma~\ref{lem:band_projector_sinc}). Consequently any framework whose
null model is a flat band reproduces the sine kernel automatically, and
the appearance of sine-kernel behaviour in such a framework is a
restatement of the null rather than evidence about the zeros of $\zeta$.
This applies in particular to the mean square of $M(x)/\sqrt{x}$ in
logarithmic average and to the energy deficit derived from it; by
Corollary~\ref{cor:profile_blindness} both are reflection-invariant and
therefore contain no cross-excursion two-point mass at all
(Remark~\ref{rem:S2_excluded}).

The criterion that separates the informative case is the following.

\begin{quote}
A statistic exhibits sine-kernel content non-tautologically if and only
if its value under the null is a number different from its value under
the sine kernel. The measured position between the two numbers is then
the sine-kernel coefficient.
\end{quote}

Three quantities measured in separate work satisfy this criterion, and
for each of them the sine kernel is written in
Section~\ref{sec:sine_kernel_from_measurements} as an explicit function
of the measurement: the consecutive spacing ratio through the
superposition family \eqref{eq:superposition_cluster}, the
freeness-defect sequence through the triangular moments
\eqref{eq:triangle_moments}, and the lobe spacing variance through the
anchors \eqref{eq:spacing_variance_anchors}.

\section{Standing results used as inputs}
\label{sec:standing_results}

The following results are established in separate work by the author.
They are stated here without proof and are used, where indicated, as
inputs. No claim of novelty is made for them in this notebook, and no
argument below reproves them.

\begin{remark}
\label{rem:standing_results_list}
\emph{Unconditional results.}

\begin{enumerate}[label=(S\arabic*)]

\item \label{item:S1}
Two-sided growth bounds for the antisymmetrised operator,
\begin{equation}
c\sqrt{N} \;\le\; \lambda_{\max}(D_N) \;\le\; C_\varepsilon\,
N^{1/2+\varepsilon},
\label{eq:DN_growth}
\end{equation}
together with the associated Cauchy interlacing corollaries.
\emph{Used as context in
Remark~\ref{rem:negative_result_density_of_states}.}

\item \label{item:S2}
Trace-norm convergence of $G_N/N^2$ to a compact limit operator.
\emph{Context.}

\item \label{item:S3}
The exact Chinese-remainder formula for the entries of $Q_N$.
\emph{Context; it is the source of \ref{item:S4}.}

\item \label{item:S4}
The tensor factorisation
\begin{equation}
\widehat{K}(r,q) \;=\; \prod_{p} L_p\!\left( v_p(r), v_p(q)\right)
\label{eq:tensor_factorisation_standing}
\end{equation}
over all of $\mathbb{N} \times \mathbb{N}$, with explicit tridiagonal
local factors $L_p$, and with the sign error in the corresponding
proposition corrected. \emph{Used actively: this is hypothesis (G2) of
Definition~\ref{def:hypotheses_on_G}.}

\item \label{item:S5}
Positivity of the Hermitian collision coefficient: for every nontrivial
zero $\rho = \beta + i\gamma$,
\begin{equation}
C_{\mathrm{herm}}(\rho)
= \frac{2\left(\gamma^2 - \beta^2\right)}
{|\rho|^4\,|\zeta'(\rho)|^2} \;>\; 0,
\label{eq:C_herm}
\end{equation}
established without assuming the Riemann hypothesis. \emph{Context; see
Remark~\ref{rem:extremal_abscissa}.}

\item \label{item:S6}
The closed form of the second-order curvature of the twisted product at
$\vartheta = 0$: for squarefree $n$ with $r$ prime factors,
\begin{equation}
z_0''(n) = -\frac{r^2}{2}\,\Var_{p \mid n}\!\left(\log p\right),
\label{eq:z0_second_derivative}
\end{equation}
together with its Dirichlet series
\begin{equation}
K(s) = -\sum_{n} \frac{z_0''(n)}{n^{s}}
= \frac{\zeta(s)}{2}\left( S_0 S_2 - S_1^2 \right),
\qquad S_2 = \left(\zeta'/\zeta\right)'.
\label{eq:K_of_s}
\end{equation}
\emph{Context.}

\item \label{item:S7}
The closed form of $z_\vartheta(n)$ for squarefree $n$ with $r$ prime
factors, as a finite cosine sum over ordered partitions of the prime
set via the cocycle $\Phi(\pi)$, with exact fourth-order coefficients
\begin{equation}
A_r = \frac{3(6r-11)}{8},
\qquad
B_r = \frac{r(6r-11)}{8},
\label{eq:fourth_order_coefficients}
\end{equation}
valid for every $r$. \emph{Context; it is the source of the twisted
action considered in Conjecture~\ref{conj:twisted_orientation_bits}.}

\item \label{item:S8}
Blindness to $\Theta$ of the entire tower of fixed transversal jets,
with complete proof, via the critical root of
$P_x(z) = \sum_n \mu(n)^2 z^{\Omega(n)}$ near $z = -1$, the limit
$\Phi_x(u) \to u e^{-u}$ uniform on compacta, and Cauchy bounds.
\emph{Context; the restriction to \emph{fixed} jets is what leaves the
scaling window of
Conjecture~\ref{conj:twisted_orientation_bits} open.}

\item \label{item:S9}
The unconditional equivalence between the Riemann hypothesis and tail
stability of a two-channel positive logarithmic deformation built from
the ordinal-prime transform. \emph{Context; see
Question~\ref{q:robin_cone_comparison}.}

\item \label{item:S10}
The exact permutation-null mean
\begin{equation}
\frac{\mathbb{E}_\pi[E_N]}{N}
= \frac{1}{\pi^2} + \frac{M(N)^2}{3N} + o(1),
\qquad E_N = \frac{1}{N}\sum_{n \le N} M(n)^2 .
\label{eq:permutation_null_mean}
\end{equation}
\emph{Used actively in Section~\ref{sec:three_conditionings}.}

\end{enumerate}
\end{remark}

\begin{remark}
\label{rem:extremal_abscissa}
One consequence of \ref{item:S5} deserves recording, since it narrows a
gap that was previously stated more broadly. At the extremal abscissa
$s = 2\Theta$ every contribution arises from zeros with $\beta = \Theta$
and carries the coefficient \eqref{eq:C_herm}, which is positive. A sum
of positive terms does not cancel, so cancellation between distinct
pairs cannot occur \emph{at the extremal abscissa}. What remains open is
the case in which $\Theta$ is a supremum that is not attained, where a
limiting or Tauberian formulation is required. We do not treat it here.
\end{remark}

\section{The three-constraint model}
\label{sec:three_conditionings}

\subsection{Mask, bits, and the fixed-count bridge}
\label{subsec:mask_and_bits}

\begin{definition}
\label{def:mask}
Fix $N \ge 2$. The \emph{squarefree mask} at level $N$ is
\begin{equation}
\mathcal{M}_N = \{\, n \le N : \mu(n) \neq 0 \,\},
\qquad
m = m(N) = |\mathcal{M}_N| ,
\label{eq:mask_definition}
\end{equation}
and the \emph{orientation bits} of the M\"obius path are
\begin{equation}
\varepsilon = (\varepsilon_n)_{n \in \mathcal{M}_N} \in \{\pm 1\}^{m},
\qquad
\varepsilon_n = \mu(n).
\label{eq:orientation_bits}
\end{equation}
\end{definition}

\begin{definition}
\label{def:fixed_count_bridge}
Let $m_{\pm}$ be the number of $n \in \mathcal{M}_N$ with
$\mu(n) = \pm 1$, so $m_+ + m_- = m$ and $m_+ - m_- = M(N)$. The
\emph{fixed-count sign bridge} $\pi$ is the uniform measure on the
$\binom{m}{m_+}$ vectors of $\{\pm 1\}^m$ with $m_+$ entries equal to
$+1$. Write
\begin{equation}
\bar{\varepsilon} = \frac{M(N)}{m}.
\label{eq:mean_bit}
\end{equation}
\end{definition}

Definition~\ref{def:fixed_count_bridge} is ingredient \textbf{(N2)}:
mask frozen, sign count frozen, only the arrangement randomised.

\subsection{Exact second moment of the bridge}
\label{subsec:bridge_second_moment}

\begin{proposition}
\label{prop:bridge_covariance}
Under $\pi$,
\begin{equation}
\mathbb{E}_\pi\!\left[\varepsilon \varepsilon^{\mathsf{T}}\right]
= (1-\beta)\, I + \beta\, J,
\qquad
\beta = \frac{m \bar{\varepsilon}^{\,2} - 1}{m-1},
\label{eq:bridge_covariance}
\end{equation}
with $I$ the identity and $J$ the all-ones matrix of size $m$.
\end{proposition}

\begin{proof}
The diagonal entries are $\varepsilon_n^2 = 1$. For $n \neq n'$,
counting ordered pairs of distinct positions,
\begin{equation}
\mathbb{E}_\pi[\varepsilon_n \varepsilon_{n'}]
= \frac{m_+(m_+-1) + m_-(m_--1) - 2 m_+ m_-}{m(m-1)}
= \frac{(m_+-m_-)^2 - m}{m(m-1)} .
\label{eq:offdiagonal_computation}
\end{equation}
Since $m_+ - m_- = m\bar{\varepsilon}$, the numerator is
$m(m\bar{\varepsilon}^{\,2}-1)$.
\end{proof}

\begin{corollary}
\label{cor:bridge_projection_mean}
For an orthogonal projection $P$ on $\mathbb{R}^m$ of rank $d$,
\begin{equation}
\mathbb{E}_\pi\!\left[\langle \varepsilon, P \varepsilon\rangle\right]
= (1-\beta)\, d + \beta\, \|P\mathbf{1}\|^2 ,
\label{eq:bridge_projection_mean}
\end{equation}
and if $M(N) = o(\sqrt{m})$ then $\beta = O(1/m)$ and
\begin{equation}
\mathbb{E}_\pi\!\left[\langle \varepsilon, P \varepsilon\rangle\right]
= d + O(1).
\label{eq:bridge_projection_mean_asymptotic}
\end{equation}
\end{corollary}

\begin{proof}
Trace \eqref{eq:bridge_covariance} against $P$, using $\tr P = d$ and
$\tr(PJ) = \|P\mathbf{1}\|^2$. For the second statement,
$|\beta| \le 1/(m-1)$ once $M(N)^2 \le m$, and
$\|P\mathbf{1}\|^2 \le m$.
\end{proof}

\begin{remark}
\label{rem:N1_usage}
Ingredient \textbf{(N1)}, the factorisation of
$\Cov_\pi\!\left(M(n)^2, M(n')^2\right)$ through the three-block
decomposition determined by $(n,\, n'-n,\, N-n')$ and the resulting
$O(m)$ evaluation of the variance, is recorded in separate work and is
not reproduced here. Only its structural consequence is used, namely
that the covariance kernel is separable of bounded rank, so that the
variance of any linear statistic of the prefix squares is a finite
combination of power sums of the mask counting function. By
Proposition~\ref{prop:exact_isotropy} the third conditioning replaces
this by an exact closed form, so \textbf{(N1)} is needed only at the
level of the bridge.
\end{remark}

\section{The reflection fibre of the quadratic profile}
\label{sec:reflection_fibre}

This section and the next develop ingredient \textbf{(N3)}.

\begin{definition}
\label{def:quadratic_profile}
The \emph{quadratic profile} at level $N$ is
\begin{equation}
\Pi_N(\varepsilon) = \left( M(n)^2 \right)_{n \le N}.
\label{eq:quadratic_profile}
\end{equation}
\end{definition}

\begin{definition}
\label{def:excursions}
Let $0 = z_0 < z_1 < \dots < z_{Z}$ enumerate
$\{ n \le N : M(n) = 0 \} \cup \{0\}$, with $Z = Z(N)$. The
\emph{excursions} are
\begin{equation}
E_j = \{\, n : z_{j-1} < n \le z_j \,\},
\qquad j = 1,\dots,Z,
\label{eq:excursion_blocks}
\end{equation}
together with a terminal block
$E_{Z+1} = \{ n : z_Z < n \le N \}$, empty exactly when $M(N) = 0$.
\end{definition}

\begin{definition}
\label{def:reflection_group}
The \emph{excursion reflection group} is
$\mathcal{R}_N \cong (\mathbb{Z}/2)^{\,Z(N)}$, acting by
\begin{equation}
(\sigma \cdot \varepsilon)_n = \sigma_{j(n)}\, \varepsilon_n,
\qquad \sigma \in \{\pm 1\}^{Z},
\label{eq:reflection_action}
\end{equation}
where $j(n)$ indexes the excursion containing $n$ and the terminal
block is held fixed.
\end{definition}

\begin{proposition}
\label{prop:reflection_fibre}
For every $\sigma \in \mathcal{R}_N$ and every $n \le N$,
\begin{equation}
M_{\sigma \cdot \varepsilon}(n) = \sigma_{j(n)}\, M_\varepsilon(n),
\label{eq:path_equivariance}
\end{equation}
so $\Pi_N$ is $\mathcal{R}_N$-invariant. Conversely, two sign vectors
on the same mask with the same quadratic profile differ by an element
of $\mathcal{R}_N$. Hence the fibre of $\Pi_N$ through $\varepsilon$ is
exactly the orbit $\mathcal{R}_N \cdot \varepsilon$.
\end{proposition}

\begin{proof}
For $n \in E_j$ we have $M_\varepsilon(z_{j-1}) = 0$, hence
$M_\varepsilon(n) = \sum_{z_{j-1} < k \le n} \varepsilon_k$ is a sum
entirely inside $E_j$; applying $\sigma$ multiplies every summand by
$\sigma_j$, giving \eqref{eq:path_equivariance}.

Conversely, let $\varepsilon, \varepsilon'$ have the same mask and
satisfy $M_{\varepsilon}(n)^2 = M_{\varepsilon'}(n)^2$ for all $n$.
The two paths then have the same zero set, hence the same excursion
decomposition, and $|M_\varepsilon| = |M_{\varepsilon'}|$ pointwise. On
each $E_j$ the path does not vanish in the interior, so $M_\varepsilon$
has constant sign there; let $\sigma_j$ be the ratio of the two signs.
A path on a fixed mask is determined by its increments, and the
increments of $M_{\varepsilon'}$ on $E_j$ are $\sigma_j$ times those of
$M_\varepsilon$, so $\varepsilon' = \sigma \cdot \varepsilon$.
\end{proof}

\subsection{The canonical null on orientation bits}
\label{subsec:canonical_null}

\begin{definition}
\label{def:excursion_signs}
The \emph{excursion signs} of $\varepsilon$ are
$s = (s_1,\dots,s_Z) \in \{\pm 1\}^Z$, where $s_j$ is the sign of
$M_\varepsilon$ in the interior of $E_j$. The \emph{reflection null}
$\nu$ is the uniform measure on $\{\pm 1\}^Z$.
\end{definition}

\begin{proposition}
\label{prop:exact_isotropy}
Under $\nu$,
\begin{equation}
\mathbb{E}_\nu\!\left[s s^{\mathsf{T}}\right] = I_Z
\label{eq:exact_isotropy}
\end{equation}
exactly, with no rank-one correction. Hence for symmetric $A$,
\begin{equation}
\mathbb{E}_\nu\!\left[\langle s, A s\rangle\right] = \tr A,
\qquad
\Var_\nu\!\left[\langle s, A s\rangle\right]
= 2 \sum_{j \neq k} A_{jk}^2 ,
\label{eq:exact_mean_variance}
\end{equation}
and for an orthogonal projection $\widehat{P}$ of rank $d$,
\begin{equation}
\mathbb{E}_\nu\!\left[\langle s, \widehat{P} s\rangle\right] = d,
\qquad
\Var_\nu\!\left[\langle s, \widehat{P} s\rangle\right]
= 2\left( d - \sum_{j} \widehat{P}_{jj}^2 \right).
\label{eq:exact_projection_moments}
\end{equation}
\end{proposition}

\begin{proof}
The $s_j$ are independent Rademacher variables under $\nu$, giving
\eqref{eq:exact_isotropy}. Writing
$\langle s, As\rangle = \tr A + \sum_{j \neq k} A_{jk} s_j s_k$ and
using $\mathbb{E}_\nu[s_j s_k s_{j'} s_{k'}] = 1$ exactly when
$\{j,k\} = \{j',k'\}$ gives the variance. For a projection,
$\sum_{j,k}\widehat{P}_{jk}^2 = \tr \widehat{P}^2 = d$.
\end{proof}

\begin{remark}
\label{rem:hierarchy_of_nulls}
Corollary~\ref{cor:bridge_projection_mean} and
Proposition~\ref{prop:exact_isotropy} form a hierarchy. Conditioning on
the mask removes the squarefree density, the source of the term
$1/\pi^2$ in \eqref{eq:permutation_null_mean}; conditioning on the sign
count removes the endpoint drift $M(N)^2/(3N)$; conditioning in
addition on the quadratic profile leaves exactly $\mathcal{R}_N$, whose
null is isotropic \emph{without} the $O(1)$ correction of
\eqref{eq:bridge_projection_mean_asymptotic}. The third conditioning is
the step at which the null becomes parameter-free.
\end{remark}

\section{The divisor-defect Gram form on orientation bits}
\label{sec:gram_form}

\begin{definition}
\label{def:divisor_defect_form}
Let $B$ be the divisor-defect matrix attached to $\mathcal{M}_N$ and
\begin{equation}
G = B^{\mathsf{T}} B
\label{eq:gram_form}
\end{equation}
its Gram form on $\mathbb{R}^{m}$.
\end{definition}

\begin{definition}
\label{def:hypotheses_on_G}
$G$ is \emph{admissible} if:
\begin{enumerate}[label=(G\arabic*)]
\item $G$ is real symmetric and positive semidefinite;
\item the associated kernel admits the tensor factorisation
\eqref{eq:tensor_factorisation_standing} of \ref{item:S4}, with local
tridiagonal factors indexed by prime valuations.
\end{enumerate}
\end{definition}

Property (G2) is \ref{item:S4} and is not proved here. Its role is to
make the eigenbasis of $G$ a character basis of the multiplicative
structure; see Remark~\ref{rem:multiplicative_band}.

\begin{definition}
\label{def:pushed_form}
The \emph{orientation-bit form} is
\begin{equation}
\widehat{G}_{jk} \;=\; \sum_{a \in E_j} \sum_{b \in E_k} G_{ab},
\qquad 1 \le j,k \le Z,
\label{eq:pushed_form}
\end{equation}
with spectral resolution
\begin{equation}
\widehat{G} = \sum_{k} \lambda_k \widehat{P}_k,
\qquad d_k = \rank \widehat{P}_k .
\label{eq:pushed_spectral}
\end{equation}
\end{definition}

\begin{proposition}
\label{prop:pushed_form_odd}
The form $s \mapsto \langle s, \widehat{G} s\rangle$ is invariant under
the global flip $s \mapsto -s$, and its off-diagonal part is odd under
every individual reflection $s_j \mapsto -s_j$. In particular
$\langle s, \widehat{G} s\rangle$ is not a function of the quadratic
profile unless $\widehat{G}$ is diagonal.
\end{proposition}

\begin{proof}
The term $s_j s_k$ changes sign under $s_j \mapsto -s_j$ for
$j \neq k$, while $\widehat{G}_{jj} s_j^2 = \widehat{G}_{jj}$ is
constant. By Proposition~\ref{prop:reflection_fibre} any function of
the profile is $\mathcal{R}_N$-invariant.
\end{proof}

\section{Blindness of the quadratic profile}
\label{sec:blindness}

\begin{lemma}
\label{lem:energy_as_two_point}
For every $N \ge 1$,
\begin{equation}
\sum_{n \le N} M(n)^2
= \sum_{a,b \le N} \mu(a)\mu(b) \left( N + 1 - \max(a,b) \right).
\label{eq:energy_as_two_point}
\end{equation}
\end{lemma}

\begin{proof}
Expand $M(n)^2 = \sum_{a,b \le n} \mu(a)\mu(b)$ and exchange the order
of summation; the pair $(a,b)$ occurs for exactly $N+1-\max(a,b)$
values of $n$.
\end{proof}

\begin{theorem}
\label{thm:cross_excursion_vanishing}
Assume $M(N) = 0$, so the excursion decomposition has no terminal
block. Then for all $j \neq k$,
\begin{equation}
\sum_{a \in E_j} \sum_{b \in E_k}
\mu(a)\mu(b)\left( N + 1 - \max(a,b) \right) = 0 .
\label{eq:cross_excursion_vanishing}
\end{equation}
Consequently $\sum_{n \le N} M(n)^2$ depends only on the
within-excursion two-point data.
\end{theorem}

\begin{proof}
Let $j < k$, so every $a \in E_j$ precedes every $b \in E_k$ and
$\max(a,b) = b$. The double sum factorises as
\begin{equation}
\left( \sum_{a \in E_j} \mu(a) \right)
\left( \sum_{b \in E_k} \mu(b)\left(N+1-b\right) \right),
\label{eq:cross_excursion_factorisation}
\end{equation}
and $\sum_{a \in E_j} \mu(a) = M(z_j) - M(z_{j-1}) = 0$ by
Definition~\ref{def:excursions}. The case $j > k$ is symmetric.
\end{proof}

\begin{corollary}
\label{cor:profile_blindness}
No statistic factoring through $\Pi_N$ detects cross-excursion
correlation of $\mu$. In particular the energy, all of its moments, and
all functionals of excursion heights and lengths are blind to it.
\end{corollary}

\begin{proof}
A statistic factoring through $\Pi_N$ is $\mathcal{R}_N$-invariant by
Proposition~\ref{prop:reflection_fibre}; by
Theorem~\ref{thm:cross_excursion_vanishing} the cross-excursion
two-point mass contributes nothing to the energy, and by
\eqref{eq:reflection_action} it is exactly the part of the two-point
data on which $\mathcal{R}_N$ acts nontrivially.
\end{proof}

\begin{remark}
\label{rem:blindness_interpretation}
The vanishing in \eqref{eq:cross_excursion_vanishing} is caused by the
defining property of an excursion, that its increments sum to zero, and
therefore cannot be repaired by reweighting: any weight depending on
$a$ and $b$ only through $\max(a,b)$ yields the same factorisation
\eqref{eq:cross_excursion_factorisation}.
\end{remark}

\begin{remark}
\label{rem:terminal_block}
If $M(N) \neq 0$ the terminal block is nonempty and
\eqref{eq:cross_excursion_vanishing} fails for pairs involving it, since
$\sum_{a \in E_{Z+1}} \mu(a) = M(N)$; the correction is a single
rank-one term proportional to $M(N)$. We normalise by taking $N$ at a
return point.
\end{remark}

\section{The scale-resolved statistic}
\label{sec:scale_resolved}

\begin{definition}
\label{def:scale_resolved_statistic}
With \eqref{eq:pushed_spectral} and $s$ the excursion signs of the true
M\"obius path,
\begin{equation}
\mathcal{E}_k = \langle s, \widehat{P}_k s \rangle,
\qquad
\mathcal{R}_k = \frac{\mathcal{E}_k}{d_k}.
\label{eq:scale_resolved_statistic}
\end{equation}
\end{definition}

\begin{proposition}
\label{prop:statistic_calibration}
Under $\nu$,
\begin{equation}
\mathbb{E}_\nu\!\left[\mathcal{R}_k\right] = 1
\quad\text{exactly},
\qquad
\Var_\nu\!\left[\mathcal{R}_k\right]
= \frac{2}{d_k^2}\left( d_k - \sum_j \widehat{P}_{k,jj}^2 \right),
\label{eq:statistic_calibration}
\end{equation}
and $\sum_k \mathcal{E}_k = \|s\|^2 = Z = \sum_k d_k$.
\end{proposition}

\begin{proof}
Apply Proposition~\ref{prop:exact_isotropy} with
$A = \widehat{P}_k$ and use $\sum_k \widehat{P}_k = I_Z$.
\end{proof}

\begin{definition}
\label{def:anisotropy}
The \emph{Gram anisotropy} at level $N$ is $k \mapsto \mathcal{R}_k-1$.
\end{definition}

\begin{remark}
\label{rem:no_free_parameters}
The calibration \eqref{eq:statistic_calibration} has no free parameter,
no fitted constant and no asymptotic normalisation. By
Proposition~\ref{prop:pushed_form_odd} and
Corollary~\ref{cor:profile_blindness}, $\mathcal{R}_k$ is
reflection-odd and therefore measures exactly the information the
quadratic profile cannot carry.
\end{remark}

\begin{conjecture}
\label{conj:anisotropy_nontrivial}
The Gram anisotropy is not identically $o(1)$: there are eigenspaces of
$\widehat{G}$ on which $\mathcal{R}_k - 1$ exceeds the null standard
deviation \eqref{eq:statistic_calibration} by an unbounded factor as
$N \to \infty$.
\end{conjecture}

\section{Sine-kernel content from measured quantities}
\label{sec:sine_kernel_from_measurements}

\subsection{Flat bands, and why they make the question empty}
\label{subsec:flat_bands}

\begin{lemma}
\label{lem:band_projector_sinc}
Let $P$ be the orthogonal projection onto a contiguous block of $d$
characters of a cyclic group of order $Z$. In the position basis the
kernel of $P$ is a Dirichlet kernel, and under the scaling
$Z \to \infty$ with $d/Z$ fixed and positions rescaled by the mean
spacing it converges to $\sin(\pi x)/(\pi x)$.
\end{lemma}

\begin{proof}
$P(n,n') = Z^{-1}\sum_{t \in S} e^{2\pi i t(n-n')/Z}$ with $S$ an
interval of length $d$ is a Dirichlet kernel, whose scaling limit is as
stated; see \cite{Mehta2004}.
\end{proof}

Equivalently, on the Fourier side the unitary two-level form factor
$\min(|\alpha|,1)$ is the indicator of a band, and
\begin{equation}
1 - R_2(x) \;=\; Y_2(x) \;=\;
\left(\frac{\sin \pi x}{\pi x}\right)^{2}
\;=\; \int_{-1}^{1}\left(1-|\tau|\right) e^{2\pi i \tau x}\, d\tau ,
\label{eq:fejer_identity}
\end{equation}
the Fej\'er identity: the sine-kernel cluster function is the transform
of the triangle. Montgomery's result gives the pair correlation for
$|\alpha| < 1$ \cite{Montgomery1973}; the plateau $|\alpha| > 1$ is
open.

\begin{warning}
\label{warn:sinc_is_the_null}
By Proposition~\ref{prop:exact_isotropy} the reflection null has
$\mathbb{E}_\nu[s s^{\mathsf{T}}] = I$, which restricted to a
contiguous block of characters is, by
Lemma~\ref{lem:band_projector_sinc}, a sine kernel. Hence the sine
kernel is the \emph{null hypothesis} of $\mathcal{R}_k$, not a
consequence of it. Deriving sine-kernel behaviour from flatness of the
null is deriving the null from the null and carries no evidential
weight concerning the zeros of $\zeta$.
\end{warning}

\begin{remark}
\label{rem:S2_excluded}
The same exclusion applies to two quantities that might appear to be
candidates. The logarithmically averaged mean square of
$M(x)/\sqrt{x}$ tends to $\sum_\rho |\rho \zeta'(\rho)|^{-2}$, a
diagonal quantity; and the ratio of the true quadratic energy to
\eqref{eq:permutation_null_mean} is a function of that same quantity.
Both are $\mathcal{R}_N$-invariant, so by
Corollary~\ref{cor:profile_blindness} neither contains any
cross-excursion two-point mass. Neither admits a nontrivial
sine-kernel reading, since neither sees $Y_2$ away from the origin.
\end{remark}

\subsection{The criterion}
\label{subsec:criterion}

\begin{remark}
\label{rem:nontautological_criterion}
A statistic carries sine-kernel content non-tautologically if and only
if its value under the null is a number different from its value under
the sine kernel. The measured position between the two numbers is then
the sine-kernel coefficient. Warning~\ref{warn:sinc_is_the_null}
describes exactly the failure of this criterion; the three subsections
below give quantities that satisfy it.
\end{remark}

\subsection{Through the consecutive spacing ratio}
\label{subsec:spacing_ratio}

\begin{definition}
\label{def:superposition}
For $p \ge 1$, the \emph{superposition family} consists of the
superposition of $p$ independent identically distributed spectra, each
unfolded to density $1/p$. Its two-level cluster function is
\begin{equation}
Y_2^{(p)}(x) = \frac{1}{p}\, Y_2\!\left(\frac{x}{p}\right),
\label{eq:superposition_cluster}
\end{equation}
so that with $Y_2 = \left(\sin \pi x/\pi x\right)^2$,
\begin{equation}
R_2^{(p)}(x) = 1 - \frac{1}{p}
\left(\frac{\sin(\pi x/p)}{\pi x/p}\right)^{2},
\qquad
Y^{(p)}_{jk} = \frac{1}{p}
\left(\frac{\sin\!\left(\pi (j-k)/p\right)}{\pi (j-k)/p}\right)^{2}.
\label{eq:superposition_matrix}
\end{equation}
\end{definition}

The coefficient $1/p$ is the \emph{sine-kernel fraction}; $p = 1$ is
the pure sine kernel and $p \to \infty$ is the vanishing of local
repulsion.

\begin{remark}
\label{rem:superposition_sum_rule}
The identity $\int Y_2^{(p)} = 1$ holds for every finite $p$, whereas a
Poisson spectrum has $\int Y_2 = 0$; the limit $p \to \infty$ is
therefore non-uniform between local behaviour and the sum rule. This
does not affect the local reading used here.
\end{remark}

\begin{experiment}
\label{exp:spacing_ratio}
\emph{Recorded in separate work, not performed here.} The consecutive
spacing ratio statistic $\tilde{r}$ of the spectrum of the smooth
arithmetic term $2N/[m,n]$ was measured and found to be approximately
$0.42$.
\end{experiment}

\begin{remark}
\label{rem:spacing_ratio_anchors}
The reference values of $\tilde{r}$ are $2\ln 2 - 1 = 0.38629\ldots$
for a Poisson spectrum, approximately $0.5307$ for the orthogonal class
and approximately $0.5996$ for the unitary class
\cite{AtasEtAl2013}; the last two are recalled and have not been
verified here.
\end{remark}

\begin{proposition}
\label{prop:positive_sine_fraction}
Suppose $\tilde{r}$ is strictly monotone in $p$ along the family of
Definition~\ref{def:superposition}, with $\tilde{r}(1)$ the unitary
value and $\tilde{r}(\infty) = 2\ln 2 - 1$. If a measured value
satisfies $\tilde{r} > 2\ln 2 - 1$ strictly, then the corresponding
$p$ is finite and the sine-kernel fraction $1/p$ in
\eqref{eq:superposition_matrix} is strictly positive.
\end{proposition}

\begin{proof}
Immediate from strict monotonicity and the stated endpoint values.
\end{proof}

\begin{remark}
\label{rem:sine_fraction_is_measured}
Proposition~\ref{prop:positive_sine_fraction} is the point at which the
criterion of Remark~\ref{rem:nontautological_criterion} is met. A flat
null forces $\tilde{r}$ to the single value $2\ln 2 - 1$ and can
produce nothing else, whereas the value recorded in
Experiment~\ref{exp:spacing_ratio} exceeds it. The sine-kernel
coefficient is then not assumed but determined, by solving
$\tilde{r}(p) = 0.42$ for $p$ and substituting into
\eqref{eq:superposition_matrix}. The monotonicity hypothesis of
Proposition~\ref{prop:positive_sine_fraction} is standard for
superpositions but is stated as a hypothesis, not proved here.
\end{remark}

\subsection{Through the moments of the triangular deviation}
\label{subsec:triangle_moments}

By \eqref{eq:fejer_identity}, sine-kernel behaviour is equivalent to the
two-level form factor deviating from its flat value by the triangle,
\begin{equation}
K(\tau) - 1 = -\left(1-\tau\right)_{+} .
\label{eq:form_factor_deviation}
\end{equation}

\begin{lemma}
\label{lem:triangle_moments}
The moments of the triangle are
\begin{equation}
m_n = \int_0^1 (1-\tau)\,\tau^{n}\, d\tau
= \frac{1}{n+1} - \frac{1}{n+2}
= \frac{1}{(n+1)(n+2)} .
\label{eq:triangle_moments}
\end{equation}
Since $(1-\tau)_+$ is supported on the compact interval $[0,1]$, it is
uniquely determined by $\{m_n\}_{n \ge 0}$.
\end{lemma}

\begin{proof}
The integral is elementary; uniqueness is the determinacy of the
Hausdorff moment problem on a compact interval.
\end{proof}

\begin{remark}
\label{rem:moment_reconstruction}
Lemma~\ref{lem:triangle_moments} changes the logical status of a
defect sequence. A deviation in the variance of a linear statistic of
order $k$ is a moment of \eqref{eq:form_factor_deviation} against a
known weight. A sequence of such defects therefore does not merely
agree with the sine kernel: by determinacy it \emph{reconstructs} it.
Two defects do not determine the triangle, but their ratio already
discriminates.
\end{remark}

\begin{experiment}
\label{exp:freeness_defects}
\emph{Recorded in separate work, not performed here.} In a test of
asymptotic freeness of the arithmetic operator against the shift, the
relative defect from the free value was approximately $12\%$ at order
two and approximately $32\%$ at order three.
\end{experiment}

\begin{remark}
\label{rem:defect_ratio_discrimination}
The three candidate growth patterns from order two to order three are
distinct. The triangle moments \eqref{eq:triangle_moments} give
$m_3/m_2 = 12/20 = 0.6$, so they \emph{decrease}. The variance of
linear statistics in the Chebyshev basis is diagonal with variance
growing linearly in the order \cite{Johansson1998}, which would give
$3/2 = 1.5$; this is recalled and not verified here. The defects of
Experiment~\ref{exp:freeness_defects} give $32/12 \approx 2.67$. The
observed growth is faster than linear, whereas the triangle requires
decay. This discrepancy is a measured quantity and is the natural
candidate for the non-sine-kernel part. The ratio is invariant under a
global rescaling of the normalisation of ``relative defect'', so the
discrimination does not depend on that convention; the individual
moments do, and cannot be read without it.
\end{remark}

\subsection{Through the lobe spacing variance}
\label{subsec:spacing_variance}

\begin{definition}
\label{def:lobe_energy}
Let $s_n$ denote the normalised spacing associated with the $n$-th
lobe, so that $s_n = 1$ characterises the ideal lobe. The
\emph{lobe curvature energy} in a window of ordinals is
\begin{equation}
\mathcal{E}_{\mathrm{lobe}}(n_0, L)
= \frac{1}{L}\sum_{n_0 < n \le n_0 + L} \left(s_n - 1\right)^2 .
\label{eq:lobe_energy}
\end{equation}
\end{definition}

\begin{lemma}
\label{lem:spacing_variance_anchors}
Under the Wigner surmises, the variance of the normalised spacing is
\begin{equation}
\Var_{\mathrm{Poisson}}(s) = 1,
\qquad
\Var_{\mathrm{GOE}}(s) = \frac{4}{\pi} - 1,
\qquad
\Var_{\mathrm{GUE}}(s) = \frac{3\pi}{8} - 1 .
\label{eq:spacing_variance_anchors}
\end{equation}
\end{lemma}

\begin{proof}
For the surmise densities $p(s) = \frac{\pi}{2} s e^{-\pi s^2/4}$ and
$p(s) = \frac{32}{\pi^2} s^2 e^{-4 s^2/\pi}$, both normalised to
$\langle s\rangle = 1$, the ratios of Gaussian moments give
$\langle s^2\rangle = 4/\pi$ and $\langle s^2\rangle = 3\pi/8$
respectively; the Poisson case is the exponential law.
\end{proof}

\begin{remark}
\label{rem:lobe_energy_criterion}
The three values in \eqref{eq:spacing_variance_anchors} are distinct,
and the null value is $1$, so \eqref{eq:lobe_energy} satisfies the
criterion of Remark~\ref{rem:nontautological_criterion}. It is the only
observable considered in this notebook that is measured directly on
ordinates rather than on an arithmetic surrogate. The reading is
coarser than that of Section~\ref{subsec:spacing_ratio}, since the gap
distribution is not a linear functional of $Y_2$; the surmises are
approximations and not exact random matrix values.
\end{remark}

\subsection{A derived measurement}
\label{subsec:derived_measurement}

\begin{remark}
\label{rem:d_resolved_sine_fraction}
Applying Section~\ref{subsec:spacing_ratio} separately within each
divisor class yields the map
\begin{equation}
d \;\longmapsto\; \frac{1}{p(d)},
\qquad \tilde{r}\!\left(p(d)\right) = \tilde{r}_d ,
\label{eq:d_resolved_sine_fraction}
\end{equation}
the profile of sine-kernel fraction resolved by divisor class. This is
a derived measurement rather than a new one: it separates the
coefficient by arithmetic scale instead of aggregating it. The
systematic observation recorded in separate work, that the anomaly at
$d = 1$ is the weakest with stronger decorrelation for $d \ge 2$,
predicts that $1/p(1)$ is the minimum of the profile. This prediction
is falsifiable and requires no new theory.
\end{remark}

\subsection{The symmetry obstruction and the twist}
\label{subsec:symmetry_obstruction}

\begin{remark}
\label{rem:symmetry_class_obstruction}
The group $\mathcal{R}_N \cong (\mathbb{Z}/2)^Z$ and the form
$\widehat{G}$ are real. A real symmetric ensemble has $\beta = 1$
statistics, whose two-level form factor
$b(\tau) = 2\tau - \tau \log(1+2\tau)$ on $\tau \le 1$
\cite{Mehta2004} is not the indicator of a band and has no plateau; the
orthogonal two-level function contains the sine kernel but is not
$1 - (\sin \pi x/\pi x)^2$. The conjectural class for the zeros of
$\zeta$ is $\beta = 2$ \cite{Montgomery1973, Odlyzko1987}. The present
construction is therefore in the wrong symmetry class, and the
obstruction is structural.
\end{remark}

\begin{conjecture}
\label{conj:twisted_orientation_bits}
Let $U(1)^Z$ act on orientation bits by
$s_j \mapsto e^{i\vartheta_j} s_j$, extending
\eqref{eq:reflection_action} from
$(\mathbb{Z}/2)^Z \subset U(1)^Z$, and let
$\widehat{G}^{(\vartheta)}$ be the Hermitian form obtained by
transporting $G$ through the twisted action, with $\vartheta$ governed
by the cocycle of \ref{item:S7}. Then the local statistics of
$\widehat{G}^{(\vartheta)}$ interpolate between the orthogonal and
unitary classes, with crossover at $\vartheta_c \asymp 1/Z$, that is,
when the antisymmetric part of the form is comparable to the mean level
spacing.
\end{conjecture}

\begin{remark}
\label{rem:pandey_mehta}
Conjecture~\ref{conj:twisted_orientation_bits} is modelled on the known
interpolating ensembles \cite{PandeyMehta1983}. What is conjectural is
not the existence of such an interpolation but the claim that the twist
of orientation bits realises it for $\widehat{G}$ at the stated scale.
Note that \ref{item:S8} closes the tower of \emph{fixed} transversal
jets, and therefore does not cover the scaling window
$\vartheta \asymp 1/Z$.
\end{remark}

\begin{warning}
\label{warn:crossover_is_necessary_not_sufficient}
Even if Conjecture~\ref{conj:twisted_orientation_bits} holds it would
establish only that the ensemble lies in a symmetry class compatible
with the zeros of $\zeta$. That is necessary and very far from
sufficient.
\end{warning}

\begin{question}
\label{q:plateau_as_flatness}
Is there a formulation of the plateau $F(\alpha) = 1$ for $\alpha > 1$
as a flatness statement for $\widehat{G}^{(\vartheta)}$ on a band of
eigenspaces, or equivalently as flatness of \eqref{eq:lobe_energy}
across a window of ordinals? A positive answer would be a
reformulation, not a proof, but it would move the plateau into a
setting carrying a transfer structure along the excursion index.
\end{question}

\section{Route toward two-point cancellation}
\label{sec:two_point}

\begin{remark}
\label{rem:chowla_lives_in_the_odd_part}
Combining Lemma~\ref{lem:energy_as_two_point},
Theorem~\ref{thm:cross_excursion_vanishing} and
Proposition~\ref{prop:pushed_form_odd}, the two-point data of $\mu$
splits into a reflection-even part carried entirely by the quadratic
profile and entirely within excursions, and a reflection-odd part
consisting of all cross-excursion pairs. Shifted pairs
$\mu(a)\mu(a+h)$ are cross-excursion for all but a vanishing proportion
of $a$ once $h$ exceeds typical excursion lengths, so the content of
two-point cancellation lies in the reflection-odd part, which is what
$\mathcal{R}_k$ measures.
\end{remark}

\begin{conjecture}
\label{conj:cross_excursion_mass}
With
\begin{equation}
C_N(h) = \sum_{a \le N-h} \mu(a)\mu(a+h),
\label{eq:shifted_correlation}
\end{equation}
the aggregate cross-excursion mass of $\{C_N(h)\}_h$ is, in the mean,
of the order predicted by $\nu$; equivalently the Gram anisotropy is
bounded on the relevant eigenspaces.
\end{conjecture}

\begin{heuristic}
\label{heur:logarithmic_weighting}
The triangular weight in \eqref{eq:energy_as_two_point} aggregates all
shifts up to $N$ with comparable mass. Logarithmic weighting is the
regime in which two-point cancellation for multiplicative functions is
presently accessible \cite{Tao2016}, building on
\cite{MatomakiRadziwill2016}. This suggests replacing the triangular
weight by a logarithmic one before pushing to orientation bits. The
remaining obstruction is uniformity in $h$: the available theorems are
for fixed or slowly growing $h$, while the aggregate requires $h$ up to
$N$. We do not claim this gap can be closed; on this route it is the
only gap.
\end{heuristic}

\begin{question}
\label{q:odd_statistic_and_uniformity}
Does passing to $\mathcal{R}_k$ change the uniformity requirement in
Heuristic~\ref{heur:logarithmic_weighting}? Since $\widehat{G}$
aggregates pairs by excursion block rather than by shift, the relevant
uniformity is in the block index, and the two need not be equivalent.
\end{question}

\section{The Robin interface}
\label{sec:robin_interface}

Ingredient \textbf{(N4)} enters here. The form $G$ is built from the
divisor defect; Robin's criterion \cite{Robin1984} is an inequality on
$\sigma(n)$ equivalent to the Riemann hypothesis, with elementary
reformulations \cite{Lagarias2002}.

\begin{question}
\label{q:robin_cone_comparison}
Suppose the truncated functional attached to $G$ satisfies an exact
elementary Robin-type equation. Does its positivity cone coincide with
the cone associated with the tail-stability criterion \ref{item:S9}? If
the cones coincide the two criteria are reparametrisations of one
another; if they differ the pair yields a positivity family not
reducible to either.
\end{question}

\begin{warning}
\label{warn:robin_novelty}
The literature on Robin-type criteria is extensive. No claim of novelty
is made for \textbf{(N4)}, and the exact elementary form should be
checked against the literature before being described as new.
\end{warning}

\section{Negative results and failed routes}
\label{sec:negative_results}

\begin{remark}
\label{rem:negative_result_energy_scalar}
The scalar quadratic energy at a single level $N$ carries no local
information: its cross terms appear with a weight essentially constant
on the scale of the mean spacing of ordinates. The defect is not the
weight but the collapse to a single number; as a function of $\log N$
the energy is almost periodic with frequency set given by differences
of ordinates and an invertible multiplier at every real frequency. The
present framework subsumes this without the auxiliary variable.
\end{remark}

\begin{remark}
\label{rem:negative_result_height_weighting}
Within the family of statistics built from $M$ alone, the weights
attached to individual zeros decay in the ordinate, so the resulting
pair-correlation measure mixes heights. This cannot be repaired by any
operation on $\log N$, which is conjugate to differences of ordinates
and not to the ordinates themselves.
\end{remark}

\begin{remark}
\label{rem:negative_result_jacobi}
The excursion decomposition admits a classical passage to continued
fractions and hence to a Jacobi matrix \cite{Flajolet1980}, which is
real symmetric and therefore lands in the orthogonal class, as in
Remark~\ref{rem:symmetry_class_obstruction}. The coincidence of the two
obstructions, in the excursion picture and in the Gram picture,
motivates Conjecture~\ref{conj:twisted_orientation_bits}.
\end{remark}

\begin{remark}
\label{rem:negative_result_topology}
A cell-complex description of the Mertens function exists, with $M(n)$
an Euler characteristic and Betti numbers given by odd divisor sums
\cite{Bjorner2011}. It does not reach two-point cancellation, because
the complex is multiplicative while shifted correlations are additive:
there is no map relating the complex at $n$ to that at $n+h$. A complex
indexed by pairs $(n, n+h)$ with Euler characteristic $C_N(h)$ would
remove the obstruction; we are not aware of such an object and make no
claim that it exists.
\end{remark}

\begin{remark}
\label{rem:negative_result_density_of_states}
Any operator whose spectrum is the set of ordinates must reproduce the
counting function $N(T) \sim (T/2\pi)\log(T/2\pi)$. The bounds
\eqref{eq:DN_growth} of \ref{item:S1} place $\lambda_{\max}(D_N)$ at
the Wigner scale, whose integrated density is semicircular and fails
that requirement by shape. $D_N$ as it stands is therefore not a
candidate spectral realisation.
\end{remark}

\begin{remark}
\label{rem:negative_result_omega_lower_bound}
Any hierarchy of bounds on the M\"obius walk must respect the
unconditional oscillation result $M(x) = \Omega_\pm(\sqrt{x})$ with
explicit constants \cite{OdlyzkoteRiele1985}; no $\varepsilon$ in such
a hierarchy is removable.
\end{remark}

\section{Open questions and future work}
\label{sec:open_questions}

\begin{enumerate}[label=(Q\arabic*)]
\item \label{item:Q1}
Does $\widehat{G}$ of \eqref{eq:pushed_form} inherit the tensor
factorisation \eqref{eq:tensor_factorisation_standing} from $G$? The
excursion blocks are defined by the zero set of $M$ and respect no
multiplicative structure, so a clean inheritance is not expected; a
partial inheritance restricted to the local factors $L_p$ with small
$p$ would be the interesting outcome. This is a finite and elementary
question, it requires no numerical work, and everything in
Sections~\ref{sec:scale_resolved}--\ref{sec:two_point} depends on it.
It is the pivotal technical question of the notebook.
\item Establish or refute Conjecture~\ref{conj:anisotropy_nontrivial},
and determine which eigenspaces carry the anisotropy.
\item Test the prediction of
Remark~\ref{rem:d_resolved_sine_fraction} that $1/p(1)$ is the minimum
of the divisor-resolved sine-kernel fraction.
\item Fix the normalisation convention under which the defects of
Experiment~\ref{exp:freeness_defects} are moments, so that
Remark~\ref{rem:defect_ratio_discrimination} can be read
quantitatively and not only through the ratio.
\item Determine the growth of $Z(N)$. Every statement in
Sections~\ref{sec:scale_resolved}--\ref{sec:two_point} is uniform in
$Z$, but the strength of the conclusions depends on it, and no claim
is made here.
\item Close the Tauberian case of Remark~\ref{rem:extremal_abscissa}.
\item Settle Question~\ref{q:plateau_as_flatness},
Question~\ref{q:odd_statistic_and_uniformity} and
Question~\ref{q:robin_cone_comparison}, the last subject to
Warning~\ref{warn:robin_novelty}.
\end{enumerate}

\section*{Notebook History}
\label{sec:notebook_history}

Version 1.0:
Three-constraint model and exact bridge second moment; reflection fibre
of the quadratic profile and the canonical parameter-free null;
divisor-defect Gram form pushed to orientation bits; blindness theorem
for the quadratic profile; scale-resolved statistic with exact mean and
variance; exploratory routes toward sine-kernel statistics and
two-point cancellation; negative results recorded.

Version 2.0:
Abstract added. Introduction restructured to state separately what is
proved here, what is taken as input, and the methodological point.
Section~\ref{sec:standing_results} added, listing the results
established in separate work with the status and role of each, together
with the narrowing of the Hermitian-collision gap to the non-attained
supremum case in Remark~\ref{rem:extremal_abscissa}.
Section~\ref{sec:sine_kernel_from_measurements} rewritten around the
criterion of Remark~\ref{rem:nontautological_criterion} and extended
with three routes recovering the sine-kernel coefficient from measured
quantities: the superposition family and the spacing ratio, the
triangular moments and Hausdorff determinacy, and the lobe spacing
variance; the divisor-resolved sine-kernel fraction added as a derived
measurement. Question~\ref{item:Q1} promoted to pivotal status.

\begin{thebibliography}{99}

\bibitem{AtasEtAl2013}
Y.~Y.~Atas, E.~Bogomolny, O.~Giraud and G.~Roux,
\emph{Distribution of the ratio of consecutive level spacings in random
matrix ensembles},
Physical Review Letters \textbf{110} (2013), 084101.

\bibitem{Bjorner2011}
A.~Bj\"orner,
\emph{A cell complex in number theory},
arXiv:1101.5704.

\bibitem{Flajolet1980}
P.~Flajolet,
\emph{Combinatorial aspects of continued fractions},
Discrete Mathematics \textbf{32} (1980), 125--161.

\bibitem{Gonek1989}
S.~M.~Gonek,
\emph{On negative moments of the Riemann zeta-function},
Mathematika \textbf{36} (1989), 71--88.

\bibitem{Johansson1998}
K.~Johansson,
\emph{On fluctuations of eigenvalues of random Hermitian matrices},
Duke Mathematical Journal \textbf{91} (1998), 151--204.

\bibitem{Lagarias2002}
J.~C.~Lagarias,
\emph{An elementary problem equivalent to the Riemann hypothesis},
American Mathematical Monthly \textbf{109} (2002), 534--543.

\bibitem{MatomakiRadziwill2016}
K.~Matom\"aki and M.~Radziwi\l{}\l{},
\emph{Multiplicative functions in short intervals},
Annals of Mathematics \textbf{183} (2016), 1015--1056.

\bibitem{Mehta2004}
M.~L.~Mehta,
\emph{Random Matrices}, third edition,
Elsevier, Amsterdam, 2004.

\bibitem{Montgomery1973}
H.~L.~Montgomery,
\emph{The pair correlation of zeros of the zeta function},
in: Analytic Number Theory,
Proceedings of Symposia in Pure Mathematics \textbf{24},
American Mathematical Society, 1973, 181--193.

\bibitem{Odlyzko1987}
A.~M.~Odlyzko,
\emph{On the distribution of spacings between zeros of the zeta
function},
Mathematics of Computation \textbf{48} (1987), 273--308.

\bibitem{OdlyzkoteRiele1985}
A.~M.~Odlyzko and H.~J.~J.~te Riele,
\emph{Disproof of the Mertens conjecture},
Journal f\"ur die reine und angewandte Mathematik \textbf{357} (1985),
138--160.

\bibitem{PandeyMehta1983}
A.~Pandey and M.~L.~Mehta,
\emph{Gaussian ensembles of random Hermitian matrices intermediate
between orthogonal and unitary ones},
Communications in Mathematical Physics \textbf{87} (1983), 449--468.

\bibitem{Robin1984}
G.~Robin,
\emph{Grandes valeurs de la fonction somme des diviseurs et hypoth\`ese
de Riemann},
Journal de Math\'ematiques Pures et Appliqu\'ees \textbf{63} (1984),
187--213.

\bibitem{Tao2016}
T.~Tao,
\emph{The logarithmically averaged Chowla and Elliott conjectures for
two-point correlations},
Forum of Mathematics, Pi \textbf{4} (2016), e8.

\bibitem{Titchmarsh1986}
E.~C.~Titchmarsh,
\emph{The Theory of the Riemann Zeta-Function}, second edition,
revised by D.~R.~Heath-Brown,
Oxford University Press, 1986.

\end{thebibliography}

\label{sec:bibliography}

\end{document}
